\section{Dominated-set magnitude}
\label{app:magnitude}

This appendix records an additional application of the numerical algebra.
The development of the HV algorithm and its complexity proof in the main
text requires only Lebesgue measure. The magnitude viewpoint encouraged
the algebraic formulation in the earlier research report~\cite{magnitudesoftware};
it is retained here as motivation and as an extension, not as a prerequisite
for the ordinary hypervolume algorithm or as a source of a faster exponent.

\subsection{Cumulative-coordinate reduction for product measures}
Let $\nu_i$ be a nonnegative measure on the $i$th coordinate axis, finite
on the bounding interval, and let
\[
 \nu=\bigotimes_{i=1}^d\nu_i,
 \qquad F_i(t)=\nu_i([0,t]).
\]
Product structure means that a box has mass
$\nu(Q(a))=\prod_iF_i(a_i)$. It does \emph{not} mean that a union of boxes
is a Cartesian product. The following elementary reduction captures
exactly what can be factored.

\begin{proposition}[Cumulative-coordinate reduction]\label{prop:cdf}
For finite $A\subseteq\R_{\ge0}^d$, put
$F(A)=\{(F_1(a_1),\ldots,F_d(a_d)):a\in A\}$. Then
\[
                  \nu(D(A))=\HV_d(F(A)).
\]
If all required cumulative values are supplied, the reduction needs
$O(nd)$ arithmetic and coordinate operations, preserves dimension, and
does not increase the number of generators.
\end{proposition}
\begin{proof}
Every nonempty intersection of input boxes is $Q(m_J)$, where
$m_{J,i}=\min_{a\in J}a_i$. Since each $F_i$ is nondecreasing,
$F_i(m_{J,i})=\min_{a\in J}F_i(a_i)$. The mass of the intersection
therefore equals the volume of the intersection of its transformed
anchored boxes. Finite inclusion--exclusion proves the equality for the
unions. The empty input gives zero on both sides.
\end{proof}

The proof is an identity, not an exponential-time implementation:
inclusion--exclusion is used to establish the reduction; a union algorithm
computes its right side. Atoms and flat parts of the cumulative functions
are allowed. This is generally a transformation of \emph{generators},
not an ordinary invertible change of coordinates on the entire region.
When cumulative values require preprocessing, their construction cost
must be added to the HV computation.

\subsection{What dominated-set magnitude means}
\label{sec:magnitude}

\subsubsection{From similarity to metric magnitude}
For a finite metric space $X=\{x_1,\ldots,x_m\}$, form the similarity
matrix $Z_{ij}=e^{-\operatorname{dist}(x_i,x_j)}$. When $Z$ is invertible,
its magnitude is $\ind^TZ^{-1}\ind$. Equivalently, solve $Zw=\ind$ and
sum the weights. Closely spaced points have high similarity and do not
behave like entirely independent units. This is the metric analogue of
counting developed in the magnitude literature~\cite{leinstermeckes2017}.

For a compact region, an applicable continuous formulation uses a finite
signed \emph{weight measure} $W$ supported on the region, with
\begin{equation}
 \int_D e^{-\operatorname{dist}(x,y)}\,dW(y)=1
                    \quad\text{for every }x\in D.
 \label{eq:weightmeasure}
\end{equation}
In the positive-definite setting used here, existence of such a weight
measure gives $\Mag(D)=W(D)$~\cite{leinstermeckes2017}. We construct one
explicitly for anchored box unions, so no limiting discretization of the
metric space is needed for their computation. The quantity of interest is
$\Mag(D(A))$, \emph{not} the magnitude of the finite generator set $A$.

For $\operatorname{dist}(x,y)=\|x-y\|_1$, the key identity is
\begin{equation}
 e^{-\|x-y\|_1}=\prod_{i=1}^d e^{-|x_i-y_i|}.
 \label{eq:exp-product}
\end{equation}
The sum in the distance becomes a product in the kernel. This tensor
structure connects metric magnitude with the finite product algebra of
the main text. The resulting algorithm evaluates no exponentials or
similarity-matrix inverses.

For completeness, positivity and the weight-measure rule have a short
energy explanation here. The one-dimensional Laplace kernel has positive
Fourier transform $2/(1+\xi^2)$; its tensor product is positive definite.
For finite signed measures $\eta$ on $D$, write
$E(\eta)=\iint e^{-\|x-y\|_1}\,d\eta(x)d\eta(y)$.
The energy characterization of magnitude is
$\sup_{\eta\ne0}\eta(D)^2/E(\eta)$ in this setting.
If $W$ satisfies~\eqref{eq:weightmeasure}, then
$\langle\eta,W\rangle=\eta(D)$ and $E(W)=W(D)$.
Cauchy--Schwarz gives $\eta(D)^2/E(\eta)\le W(D)$, with equality
for $\eta=W$. This justifies the rule used in the proof below and
explains its relation to the finite matrix formula.

\subsubsection{A direct proof for anchored boxes}
For $a\ge0$, let
\[
 w_a=\tfrac12\bigl(\delta_0+\delta_a+\lambda|_{[0,a]}\bigr).
\]
When $a=0$, the coincident endpoint atoms give $w_0=\delta_0$.
Splitting an elementary integral at $x$ shows that, for every $x\ge0$,
\begin{equation}
 \int_{[0,a]}e^{-|x-y|}\,dw_a(y)=e^{-(x-a)_+},
 \qquad w_a([0,a])=1+a/2,
 \label{eq:interval-potential}
\end{equation}
where $t_+=\max(t,0)$. For $0\le x\le a$, the integral of the density
is $1-(e^{-x}+e^{-(a-x)})/2$; the two atoms supply the missing terms.
For $x>a$, all contributions have the common factor $e^{-(x-a)}$.

Tensor products of these interval measures give a weight measure for
$Q(a)$ and yield $\Mag(Q(a))=\prod_i(1+a_i/2)$. More is true.

\begin{theorem}[Product computing measure for dominated sets]
\label{thm:magnitude}
Let $A$ be any finite set of nonnegative corners. With
\[
                  \mu=\bigotimes_{i=1}^d
                         (\delta_0+\tfrac12\lambda),
\]
the $\ell_1$ magnitude of its dominated region satisfies
\begin{equation}
                         \Mag(D(A))=\mu(D(A)).
 \label{eq:mag-product}
\end{equation}
The magnitude of the empty region is taken to be zero.
\end{theorem}
\begin{proof}
Assume $A=\{a^1,\ldots,a^n\}$ is nonempty. For nonempty
$J\subseteq\{1,\ldots,n\}$, let $m_J$ be the coordinatewise minimum
of its corners and let $W_J=\bigotimes_i w_{m_{J,i}}$. Define
\[
             W=\sum_{\varnothing\ne J\subseteq[n]}
                       (-1)^{|J|+1}W_J.
\]
This is a finite signed measure supported on $D(A)$. Its potential at a
nonnegative point $x$ is the same signed sum of
$\exp(-\sum_i(x_i-m_{J,i})_+)$, by
\eqref{eq:exp-product}--\eqref{eq:interval-potential}.
If $x\in D(A)$, choose $k$ with $x\le a^k$. Pair every nonempty
$J$ not containing $k$ with $J\cup\{k\}$. Their potentials agree:
replacing $m_{J,i}$ by $\min(m_{J,i},a_i^k)$ leaves
$(x_i-m_{J,i})_+$ unchanged because $a_i^k\ge x_i$.
All these pairs cancel. The unpaired subset $\{k\}$ has potential one.
Thus $W$ satisfies~\eqref{eq:weightmeasure}.

Its total mass is
\[
 W(D(A))=\sum_{\varnothing\ne J\subseteq[n]}(-1)^{|J|+1}
                      \prod_i(1+m_{J,i}/2).
\]
The last product is also $\mu(Q(m_J))$. Inclusion--exclusion for the
positive measure $\mu$ now gives $W(D(A))=\mu(D(A))$.
\end{proof}

Equation~\eqref{eq:mag-product} is the dominated-set formula
in~\cite{emmerich2026}. The proof above spells out why it is legitimate
on this class, without assuming that magnitude is a valuation on arbitrary
sets. The positive \emph{computing measure} $\mu$ is not in general the
kernel weight measure $W$. For instance, the translated interval $[1,2]$
has metric magnitude $3/2$ by translation invariance, but its mass under
the fixed-anchor measure $\delta_0+\lambda/2$ is $1/2$. Consequently,
this formula does not directly identify arbitrary Klee-measure instances
with the metric magnitude of their unions.

\subsubsection{Faces, Pareto compliance, and scale}
Expanding the product measure in~\eqref{eq:mag-product} gives
\begin{equation}
 \Mag(D(A))=\sum_{I\subseteq[d]}2^{-|I|}
                       \lambda_{|I|}(\pi_I D(A)).
 \label{eq:projection}
\end{equation}
For a nonempty region the zero-dimensional term is one; for the empty
region all terms are zero. The reason projections appear is specific to
down-sets: the section obtained by setting coordinates outside $I$ to
zero coincides with the projection onto $I$.
In two dimensions the formula is
\[
 1+\tfrac12(\text{horizontal extent}+\text{vertical extent})
                      +\tfrac14(\text{area}).
\]
In higher dimensions it combines all coordinate-projection volumes,
including the full-dimensional HV. The algorithm need not enumerate
these $2^d$ projections; one product measure contains them all.

Proper inclusion of two finite dominated regions strictly increases
$\mu$, hence magnitude. To see strictness, take $p\in D(B)\setminus D(A)$.
Closedness of $D(A)$ leaves a sufficiently small lower rectangle near $p$
outside $D(A)$; use intervals on the positive coordinates of $p$ and
origin atoms on its zero coordinates. This rectangle has positive $\mu$
mass and lies in $D(B)$. If $p=0$, its atomic mass is already positive.
This explains Pareto compliance at the region level. Adding an already
dominated generator changes neither indicator. A new box of zero
full-dimensional volume can change magnitude through a lower-dimensional
face even when it leaves ordinary HV unchanged.

Magnitude uses a metric and therefore a length scale. For the distance
$\tau\|x-y\|_1$, with $\tau>0$, the computing measure becomes
$\bigotimes_i(\delta_0+\tfrac\tau2\lambda)$. Its value is a polynomial
in $\tau$, whose degree-$d$ term is $(\tau/2)^d\HV_d(A)$
(and vanishes when the full-dimensional HV is zero).
Changing objective units without correspondingly specifying the metric
changes the indicator. Product structure gives a clear algebraic
description of this choice; it does not remove the choice.

\subsection{Magnitude and hypervolume: both reductions}
\label{sec:reductions}

\begin{corollary}[Magnitude to HV]\label{cor:mag-to-hv}
For finite nonnegative corners,
\begin{equation}
 \Mag(D(A))=\HV_d\bigl(\{\ind+a/2:a\in A\}\bigr).
 \label{eq:forward}
\end{equation}
\end{corollary}
\begin{proof}
Apply Proposition~\ref{prop:cdf} with $F_i(t)=1+t/2$ and
Theorem~\ref{thm:magnitude}.
\end{proof}

The anchor on the right of~\eqref{eq:forward} remains zero. The affine
image of $[0,a_i]$ would instead be $[1,1+a_i/2]$. Thus this is not a
statement that translating a region adds volume. Transforming the
generators and rebuilding their anchored boxes includes precisely the
face contributions represented by the atoms.

\begin{proposition}[HV to magnitude]\label{prop:reverse}
For an HV input $B\subseteq\R_{\ge0}^d$, discard corners with any zero
coordinate. If none remain, the HV is zero. Otherwise put
\[
 \delta_i=\min_{b\in B}b_i>0,
 \qquad c_i(b)=2(b_i/\delta_i-1).
\]
Then
\begin{equation}
 \HV_d(B)=\left(\prod_i\delta_i\right)
                 \Mag\bigl(D(\{c(b):b\in B\})\bigr).
 \label{eq:reverse}
\end{equation}
\end{proposition}
\begin{proof}
The discarded boxes have zero $d$-dimensional volume. Scaling the
remaining $i$th coordinates by $1/\delta_i$ multiplies their HV by
$\prod_i\delta_i^{-1}$ and gives extents at least one. For these
normalized corners, $1+c_i(b)/2=b_i/\delta_i$. Apply
\eqref{eq:forward} and undo the volume scale.
\end{proof}

Both directions require one indicator evaluation and $O(nd)$ preprocessing.
They preserve dimension and input size. Therefore an arithmetic upper
bound transfers in both directions, and so does a conditional lower bound
within a computational model supporting these arithmetic transformations.
This observation alone proves no new hardness result. Bit complexity and
restrictions on integer coordinates require their own accounting.
Zero coordinates created by~\eqref{eq:reverse} must be retained in the
\emph{magnitude} instance. Discarding them would invalidate the reduction.

\begin{example}[One example in both directions]\label{ex:two-box}
Take $A=\{(2,1),(1,2)\}$. Its HV is $2+2-1=3$, whereas its magnitude is
\[
 3+3-\tfrac94=\tfrac{15}{4}
       =1+\tfrac12(2+2)+\tfrac14\cdot3.
\]
The forward reduction gives corners $(2,3/2)$ and $(3/2,2)$, with HV
$15/4$. The reverse reduction applied to the original HV instance gives
$(2,0)$ and $(0,2)$: the magnitude of these two anchored line segments
is $2+2-1=3$. Adding $(3,0)$ to the original set leaves its HV at $3$
and increases its magnitude to $17/4$.
\end{example}

\begin{figure}[tb]
\centering
\begin{tikzpicture}[scale=1.25,>=Stealth]
 \begin{scope}
 \fill[blue!12] (0,0)--(2,0)--(2,1)--(1,1)--(1,2)--(0,2)--cycle;
 \draw[very thick,blue!60!black] (0,0)--(2,0)--(2,1)--(1,1)--(1,2)--(0,2)--cycle;
 \draw[->] (-.15,0)--(2.65,0) node[right] {$x_1$};
 \draw[->] (0,-.15)--(0,2.5) node[above] {$x_2$};
 \fill (2,1) circle(1.5pt) node[right] {\small$(2,1)$};
 \fill (1,2) circle(1.5pt) node[above right] {\small$(1,2)$};
 \fill[orange!85!black] (0,0) circle(2.2pt);
 \node at (1,-.6) {$\HV=3,\quad\Mag=15/4$};
 \end{scope}
 \node at (3.9,1.1) {$a\mapsto\ind+a/2$};
 \draw[->,thick] (3,0.7)--(4.8,0.7);
 \begin{scope}[xshift=6cm]
 \fill[teal!14] (0,0)--(2,0)--(2,1.5)--(1.5,1.5)--(1.5,2)--(0,2)--cycle;
 \draw[very thick,teal!60!black] (0,0)--(2,0)--(2,1.5)--(1.5,1.5)--(1.5,2)--(0,2)--cycle;
 \draw[->] (-.15,0)--(2.65,0) node[right] {$y_1$};
 \draw[->] (0,-.15)--(0,2.5) node[above] {$y_2$};
 \fill (2,1.5) circle(1.5pt);
 \fill (1.5,2) circle(1.5pt);
 \node at (1,-.6) {$\HV=15/4$};
 \end{scope}
\end{tikzpicture}
\caption{The transformed generators define a new anchored region. The
left region's magnitude equals the right region's ordinary area. The
origin atom and lower-dimensional contributions are represented by the
enlargement of anchored extents, not by a translation of the shaded set.}
\label{fig:shift}
\end{figure}


\subsection{Using the same numerical engine}
\label{app:magnitude-software}

On the cells of Section~\ref{sec:discrete}, the initial masses are
\[
 w_i(0)=1,\qquad w_i(k)=\frac{c_{i,k}-c_{i,k-1}}2\quad(k\ge1).
\]
Replacing the Lebesgue masses by these masses changes neither the cutoff
algebra nor the recursion. Prefix sums include the original origin atom.
Cropping shifts rank indices only: it must not insert a new atom at the
new rank zero. Compression adds existing cell masses, including atoms;
it does not reconstruct a continuous density by dividing by cell volume.
The hypotheses of Theorem~\ref{thm:main} therefore apply. Alternatively,
one can transform the generators by~\eqref{eq:forward} and call an ordinary
HV solver. Neither route improves the asymptotic exponent.

For Example~\ref{ex:two-box}, the axis coordinates are $(0,1,2)$.
The HV masses are $(0,1,1)$ and the magnitude masses are $(1,1/2,1/2)$.
The only uncovered cell in the enclosing rank grid is $(2,2)$, giving
HV $4-1=3$ and magnitude $4-1/4=15/4$.

The public interface also supports the indicator directly:
\begin{lstlisting}[language=Python]
from NumericalChanHVND.numerical_chan_local_compiled import magnitude
points = [(2, 1, 1, 1), (1, 2, 1, 1)]
assert abs(magnitude(points) - 135.0 / 16.0) < 1e-10
\end{lstlisting}
Here $(15/4)(3/2)^2=135/16$. An empty input has magnitude zero, whereas
a single all-zero generator has magnitude one. Zero-extent boxes must
therefore be retained for this application. The reflection from a
minimization region to upper-corner extents is an $\ell_1$ isometry, so
it also preserves the metric magnitude of the corresponding region.

The saved validation record includes 18 checks of the two HV--magnitude
transformations. The 64-point checkpoint instance and both checkpoint
decisions were also checked for magnitude against transformed Chan $d/2$
inputs. The 1,552-box instance has magnitude $411\,863\,688.75$, again
agreeing with the transformed reference. The shared algebra tests include
boundary atoms. These are correctness checks, not evidence that magnitude
improves the running time of ordinary HV. The larger benchmark in
Section~\ref{sec:experiments} measures ordinary hypervolume only.
